% ============================================================
%  Conception de l'Architecture
%  Project: AEGIS Entropy — AI-Based Port Scan Detection
%  Updated: June 2026 (aligned with final report)
%  Approach: problem-solution traceability
% ============================================================
\documentclass[12pt,a4paper]{article}

\usepackage[utf8]{inputenc}
\usepackage[T1]{fontenc}
\usepackage[english]{babel}
\usepackage{geometry}
\usepackage{enumitem}
\usepackage{booktabs}
\usepackage{array}
\usepackage{tabularx}
\usepackage{xcolor}
\usepackage{titlesec}
\usepackage{fancyhdr}
\usepackage{hyperref}
\usepackage{longtable}

\geometry{margin=2.5cm}

% --- Colors ---
\definecolor{sectionblue}{RGB}{15, 23, 42}
\definecolor{accentcyan}{RGB}{6, 182, 212}
\definecolor{accentgreen}{RGB}{16, 185, 129}
\definecolor{accentred}{RGB}{239, 68, 68}

\titleformat{\section}
  {\large\bfseries\color{sectionblue}}
  {\thesection.}{0.5em}{}[\vspace{-0.5em}\textcolor{accentcyan}{\rule{\textwidth}{0.4pt}}]
\titleformat{\subsection}
  {\normalsize\bfseries\color{sectionblue}}
  {\thesubsection.}{0.5em}{}

\pagestyle{fancy}
\fancyhf{}
\fancyhead[L]{\small\textcolor{gray}{Architecture — AEGIS Entropy}}
\fancyhead[R]{\small\textcolor{gray}{AI Module — S2}}
\fancyfoot[C]{\small\thepage}
\renewcommand{\headrulewidth}{0.4pt}

\newcolumntype{L}[1]{>{\raggedright\arraybackslash}p{#1}}
\newcommand{\problem}{\textcolor{accentred}{\textbf{Probl\`{e}me : }}}
\newcommand{\answer}{\textcolor{accentgreen}{\textbf{R\'{e}ponse : }}}

\begin{document}

\begin{center}
  {\Large\bfseries Conception de l'Architecture}\\[0.6em]
  {\large AEGIS Entropy}\\[0.3em]
  {\large Syst\`{e}me de D\'{e}tection de Port Scan\\[0.2em]
   par Apprentissage Automatique et D\'{e}fense Active}\\[1.2em]
  \begin{tabular}{rl}
    \textbf{Module :}  & Intelligence Artificielle --- Semestre 2 \\
    \textbf{\'{E}quipe :} & 5 membres (Ing\'{e}nierie --- ENSA K\'{e}nitra) \\
    \textbf{Encadrant :}  & Pr.\ Aniss Moumen \\
    \textbf{Date :}    & Juin 2026 \\
  \end{tabular}
\end{center}

\vspace{1em}

% ============================================================
\section{Rappel de la probl\'{e}matique}

La d\'{e}tection du port scanning --- premi\`{e}re phase de reconnaissance de
la quasi-totalit\'{e} des cyberattaques --- reste un d\'{e}fi op\'{e}rationnel
pour les \'{e}quipes SOC. Les syst\`{e}mes actuels souffrent de quatre lacunes
identifi\'{e}es lors de la phase de conception du projet :

\begin{enumerate}[nosep,leftmargin=1.5em]
  \item \textbf{C\'{e}cit\'{e} face aux scans furtifs :}
        les IDS bas\'{e}s sur signatures (Snort, Suricata) appliquent des seuils
        fixes que les attaquants contournent en espa\c{c}ant leurs sondes.
  \item \textbf{Taux de faux positifs \'{e}lev\'{e} :}
        le trafic l\'{e}gitime d\'{e}clenche fr\'{e}quemment les m\^{e}mes
        signatures que les scans malveillants, noyant les analystes sous des
        milliers d'alertes par jour.
  \item \textbf{Absence de profilage comportemental :}
        les outils classiques traitent chaque paquet isol\'{e}ment, sans analyser
        les motifs temporels et statistiques du trafic par adresse source.
  \item \textbf{R\'{e}ponse manuelle et lente :}
        m\^{e}me lorsqu'une menace est d\'{e}tect\'{e}e, le blocage d\'{e}pend
        de l'intervention humaine, avec un d\'{e}lai de r\'{e}action trop long
        pour contenir une attaque en cours.
\end{enumerate}

Ces constats ont \'{e}t\'{e} valid\'{e}s par quatre entretiens semi-directifs
conduits aupr\`{e}s d'\'{e}tudiants en ing\'{e}nierie r\'{e}seaux et
cybers\'{e}curit\'{e}. Leurs retours confirment que la surveillance manuelle
est insoutenable, que les faux positifs constituent un point de douleur majeur,
et qu'un syst\`{e}me capable d'apprentissage comportemental avec r\'{e}ponse
autonome gradu\'{e}e r\'{e}pond \`{a} un besoin r\'{e}el.

% ============================================================
\section{Exigences issues de la probl\'{e}matique}

De cette probl\'{e}matique d\'{e}coulent les exigences fonctionnelles et
non-fonctionnelles que l'architecture doit satisfaire :

\begin{table}[h!]
\centering
\small
\begin{tabularx}{\textwidth}{c L{4cm} X}
\toprule
\textbf{R\'{e}f.} & \textbf{Exigence} & \textbf{Justification} \\
\midrule
E1 & D\'{e}tection temps-r\'{e}el avec F1 $\geq$ 88\% & Les IDS classiques
    produisent trop de faux positifs et manquent les scans furtifs. \\
E2 & Apprentissage comportemental & Les r\`{e}gles statiques ne capturent pas
    les motifs temporels et statistiques du trafic. \\
E3 & R\'{e}ponse autonome gradu\'{e}e & L'intervention humaine est trop lente
    ($>$ 15 secondes) ; le blocage automatique doit \^{e}tre proportionnel
    au niveau de confiance de la d\'{e}tection. \\
E4 & Tableau de bord SOC temps-r\'{e}el & Les analystes doivent visualiser
    les scanners actifs, les scores de confiance et les actions de d\'{e}fense
    en cours. \\
E5 & Explicabilit\'{e} des d\'{e}cisions & Les alertes doivent \^{e}tre
    justifi\'{e}es (features ayant contribu\'{e} \`{a} la classification)
    pour permettre la validation humaine. \\
E6 & Isolation des environnements & L'entra\^{i}nement et le d\'{e}ploiement
    doivent fonctionner dans un laboratoire VirtualBox isol\'{e}
    (192.168.100.0/24) sans connectivit\'{e} externe. \\
\bottomrule
\end{tabularx}
\caption{Exigences fonctionnelles d\'{e}riv\'{e}es de la probl\'{e}matique.}
\label{tab:exigences}
\end{table}

% ============================================================
\section{Vue d'ensemble de l'architecture}

L'architecture AEGIS Entropy est structur\'{e}e en \textbf{trois couches}
principales, compl\'{e}t\'{e}es par une couche d'observabilit\'{e} :

\begin{center}
\small
\begin{tabular}{c}
\textbf{[Capture]} $\longrightarrow$
\textbf{[D\'{e}tection ML]} $\longrightarrow$
\textbf{[R\'{e}ponse Active]} \\
\hspace{3cm} $\downarrow$ \\
\hspace{3cm} \textbf{[Observabilit\'{e} --- Dashboard SOC]}
\end{tabular}
\end{center}

\vspace{0.5em}
\noindent Le flux de donn\'{e}es complet est le suivant :

\begin{quote}
\small
\texttt{[Trafic r\'{e}seau / Nmap XML]}
$\rightarrow$
\texttt{[Capture Layer --- Scapy, BPF, 10s/60s fen\^{e}tres]}
$\rightarrow$
\texttt{[Feature Extractor --- 9 features + 2 placeholders]}
$\rightarrow$
\texttt{[Bridge --- Load scaler.pkl, load rf\_model.pkl, classify]}
$\rightarrow$
\texttt{[Severity Mapping --- LOW / MEDIUM / HIGH]}
$\rightarrow$
\texttt{[Deception : MTD + Honeypot] + [Dashboard : Socket.IO alert]}
\end{quote}

Les sections suivantes d\'{e}crivent chaque couche en tra\c{c}ant explicitement
le lien entre le probl\`{e}me identifi\'{e} au \S1 et la solution architecturale
adopt\'{e}e.

% ============================================================
\section{Couche Capture --- de l'inspection manuelle \`{a} l'agr\'{e}gation automatis\'{e}e}

\problem Les IDS classiques traitent chaque paquet isol\'{e}ment et ne
disposent pas de m\'{e}canisme d'agr\'{e}gation temporelle. Les scans furtifs
\`{a} bas d\'{e}bit passent sous les seuils parce que les r\`{e}gles statiques
ne consid\`{e}rent pas l'historique du trafic par adresse source.

\answer La couche Capture remplace l'inspection ponctuelle par une
\textbf{agr\'{e}gation par fen\^{e}tre temporelle glissante}. Deux granularit\'{e}s
sont utilis\'{e}es :

\begin{itemize}[nosep,leftmargin=1.5em]
  \item \textbf{Fen\^{e}tre de 10 secondes} pour la d\'{e}tection des scans rapides
        et agressifs.
  \item \textbf{Fen\^{e}tre de 60 secondes} pour l'identification des scans lents
        et furtifs (low-and-slow).
\end{itemize}

\noindent Le module utilise Scapy avec un filtre BPF pour capturer le trafic
TCP/UDP sur l'interface r\'{e}seau. En mode hors-ligne (d\'{e}mo), des fichiers
Nmap XML pr\'{e}-captur\'{e}s sont utilis\'{e}s comme source, garantissant la
reproductibilit\'{e} de la d\'{e}monstration. Dix profils de scan Nmap ont
\'{e}t\'{e} g\'{e}n\'{e}r\'{e}s, couvrant des dur\'{e}es de 4.66 secondes
(scan rapide) \`{a} 505 secondes (scan furtif).

L'environnement de capture est un laboratoire VirtualBox isol\'{e} :
Kali Linux (192.168.100.10) en attaquant, Ubuntu Server (192.168.100.20)
en cible, exposant SSH (port 22) et HTTP (port 80).

\bigskip
\noindent \textbf{Tra\c{c}abilit\'{e} :} cette couche r\'{e}pond aux exigences
\textbf{E1} (donn\'{e}es de qualit\'{e} pour la d\'{e}tection) et
\textbf{E6} (isolation du laboratoire).

% ============================================================
\section{Couche D\'{e}tection ML --- des signatures statiques \`{a} l'apprentissage comportemental}

\problem Les IDS bas\'{e}s sur signatures (Snort, Suricata) comparent le trafic
observ\'{e} \`{a} une base de motifs connus. Cette approche est aveugle face aux
scans polymorphes, aux attaques zero-day, et aux outils de reconnaissance
personnalis\'{e}s. De plus, elle n\'{e}cessite une mise \`{a} jour manuelle
continue des signatures.

\answer La couche D\'{e}tection remplace les r\`{e}gles statiques par
\textbf{trois mod\`{e}les d'apprentissage automatique compl\'{e}mentaires}
fonctionnant sur 11 caract\'{e}ristiques de flux :

\subsection{Pipeline de pr\'{e}traitement}

Avant classification, les flux subissent un pipeline \`{a} quatre \'{e}tapes
avec contr\^{o}le strict de non-fuite de donn\'{e}es (\textit{anti-leakage}) :

\begin{enumerate}[nosep,leftmargin=1.5em]
  \item \textbf{Nettoyage + imputation} (m\'{e}diane, robuste aux outliers).
  \item \textbf{Capping IQR} (winsorisation) --- pr\'{e}serve tous les
        enregistrements d'attaque.
  \item \textbf{Correction log1p} pour les 5 features asym\'{e}triques
        ($|\text{skew}| > 1.0$).
  \item \textbf{StandardScaler + SMOTE} --- le scaler est ajust\'{e}
        \textbf{exclusivement} sur l'ensemble d'entra\^{i}nement.
\end{enumerate}

\subsection{Les trois mod\`{e}les}

\begin{itemize}[nosep,leftmargin=1.5em]
  \item \textbf{Random Forest} (100 arbres, bagging) --- classifieur supervis\'{e}
        principal. Robuste aux donn\'{e}es tabulaires non-lin\'{e}aires, fournit
        les scores d'importance des features (explicabilit\'{e}, exigence E5).
  \item \textbf{XGBoost} (100 estimateurs, \texttt{max\_depth=6}, gradient boosting)
        --- benchmark de comparaison pour v\'{e}rifier que Random Forest n'est pas
        sous-performant par simplicit\'{e}.
  \item \textbf{Isolation Forest} (\texttt{contamination='auto'}) --- d\'{e}tecteur
        non supervis\'{e} destin\'{e} \`{a} la d\'{e}tection des scans lents et
        inconnus en production. L'\'{e}chec sur CICIDS2017 (F1 = 9.46\%) est
        document\'{e} et expliqu\'{e} par le \textbf{paradoxe de l'anomalie
        majoritaire} : PortScan = 55.48\% des donn\'{e}es.
\end{itemize}

Tous les mod\`{e}les supervis\'{e}s sont valid\'{e}s par \textbf{10-fold
stratified cross-validation}. Les r\'{e}sultats sur l'ensemble de test
(57\,294 flux) sont :

\begin{table}[h!]
\centering
\small
\begin{tabular}{lcccc}
\toprule
\textbf{Mod\`{e}le} & \textbf{F1} & \textbf{FPR} & \textbf{AUC-ROC} & \textbf{Erreurs} \\
\midrule
Random Forest & 99.92\% & 0.11\% & 0.9997 & 51 / 57\,294 \\
XGBoost       & 99.92\% & 0.11\% & 0.9999 & 49 / 57\,294 \\
Isolation Forest & 9.46\% & 53.53\% & --- & 43\,184 / 57\,294 \\
\bottomrule
\end{tabular}
\caption{Performances des mod\`{e}les sur l'ensemble de test.}
\end{table}

\subsection{Le module Bridge --- middleware d'int\'{e}gration}

Le Bridge (\texttt{bridge/bridge.py}) connecte la capture \`{a} la d\'{e}tection.
Il charge le scaler entra\^{i}n\'{e} (\texttt{scaler.pkl}), applique le
pr\'{e}traitement sur les features entrantes, charge le mod\`{e}le Random Forest
(\texttt{rf\_model.pkl}), produit une classification et un score de confiance,
et dispatche le r\'{e}sultat vers le dashboard et le sous-syst\`{e}me de d\'{e}fense.

\bigskip
\noindent \textbf{Tra\c{c}abilit\'{e} :} cette couche r\'{e}pond aux exigences
\textbf{E1} (d\'{e}tection avec F1 $\geq$ 88\%), \textbf{E2} (apprentissage
comportemental), et \textbf{E5} (explicabilit\'{e} via feature importance).

% ============================================================
\section{Couche R\'{e}ponse Active --- du blocage manuel \`{a} la d\'{e}fense autonome gradu\'{e}e}

\problem Dans un SOC classique, la d\'{e}tection d\'{e}clenche une alerte,
mais la r\'{e}ponse d\'{e}pend d'un analyste humain qui doit manuellement
bloquer l'IP source. Ce d\'{e}lai --- souvent plusieurs minutes --- permet
\`{a} l'attaquant de compl\'{e}ter sa reconnaissance avant toute contre-mesure.

\answer La couche R\'{e}ponse Active impl\'{e}mente une \textbf{d\'{e}fense
gradu\'{e}e et autonome} bas\'{e}e sur le score de confiance du mod\`{e}le ML :

\subsection{Cartographie de s\'{e}v\'{e}rit\'{e}}

\begin{table}[h!]
\centering
\begin{tabular}{l l l}
\toprule
\textbf{Confiance} & \textbf{S\'{e}v\'{e}rit\'{e}} & \textbf{Action autonome} \\
\midrule
$<$ 0.85      & LOW    & Redirection honeypot uniquement (surveillance passive) \\
0.85 -- 0.95  & MEDIUM & Redirection honeypot + blacklist IP (UFW drop) \\
$>$ 0.95      & HIGH   & Redirection + blacklist + mutation de port (MTD) \\
\bottomrule
\end{tabular}
\caption{Cartographie confiance $\rightarrow$ action.}
\end{table}

\subsection{Modules de d\'{e}fense}

\begin{itemize}[nosep,leftmargin=1.5em]
  \item \textbf{Core Deception} --- 5 profils honeypot ; leurres l\'{e}gers
        \'{e}mulant SSH, HTTP, FTP ; active le flag
        \texttt{shadow\_node\_interaction=1} en production.
  \item \textbf{Network Mutator (MTD)} --- rotation dynamique des ports de
        service via iptables ; thread d'arri\`{e}re-plan avec PRNG et p\'{e}riodes
        de gr\^{a}ce.
  \item \textbf{Traffic Shaper} --- SYN-ACK tarpit et sabotage MSS pour
        ralentir les scanners sans les bloquer.
  \item \textbf{IP Manager} --- blackhole kernel + UFW deny avec
        d\'{e}blocage automatique bas\'{e} sur TTL.
\end{itemize}

\subsection{R\'{e}sultats empiriques}

\begin{itemize}[nosep,leftmargin=1.5em]
  \item \textbf{+420\%} de temps de reconnaissance pour l'attaquant
        (scan 5$\times$ plus lent).
  \item \textbf{85\% $\rightarrow$ 20\%} de succ\`{e}s de fingerprinting OS
        (l'attaquant ne peut plus identifier le syst\`{e}me cible).
  \item \textbf{+2.3\%} de gain en accuracy ML via la boucle de retour honeypot
        (le flag \texttt{shadow\_node\_interaction} am\'{e}liore la d\'{e}tection
        en production).
\end{itemize}

\bigskip
\noindent \textbf{Tra\c{c}abilit\'{e} :} cette couche r\'{e}pond \`{a} l'exigence
\textbf{E3} (r\'{e}ponse autonome gradu\'{e}e).

% ============================================================
\section{Couche Observabilit\'{e} --- de l'alerte opaque au tableau de bord SOC}

\problem Les outils classiques g\'{e}n\`{e}rent des milliers de faux positifs
par jour sans fournir de contexte exploitable. Les analystes perdent du temps
\`{a} trier des alertes sans comprendre pourquoi elles ont \'{e}t\'{e} d\'{e}clench\'{e}es.

\answer La couche Observabilit\'{e} fournit un \textbf{tableau de bord SOC
temps-r\'{e}e}l bas\'{e} sur Flask + Socket.IO + D3.js. Trois endpoints POST
re\c{c}oivent les \'{e}v\'{e}nements en continu :

\begin{itemize}[nosep,leftmargin=1.5em]
  \item \texttt{POST /api/alert} --- \'{e}v\'{e}nements de d\'{e}tection
        (bridge $\rightarrow$ dashboard).
  \item \texttt{POST /api/block} --- demandes de blocage IP.
  \item \texttt{POST /api/honeypot\_event} --- \'{e}v\'{e}nements de d\'{e}ception
        (modules de d\'{e}fense $\rightarrow$ dashboard).
\end{itemize}

\noindent Chaque \'{e}v\'{e}nement est \'{e}mis en temps r\'{e}el vers tous les
clients connect\'{e}s via Socket.IO, et \'{e}galement sauvegard\'{e} localement
dans \texttt{detection\_logs.json} (format JSONL). Si le dashboard est hors-ligne,
le bridge \'{e}crit les \'{e}v\'{e}nements localement sans perte de donn\'{e}es.

Le dashboard fonctionne \textbf{enti\`{e}rement hors-ligne} sur Windows, sans
d\'{e}pendance Linux.

\bigskip
\noindent \textbf{Tra\c{c}abilit\'{e} :} cette couche r\'{e}pond \`{a} l'exigence
\textbf{E4} (tableau de bord temps-r\'{e}el).

% ============================================================
\section{R\'{e}sum\'{e} : tra\c{c}abilit\'{e} probl\`{e}me $\rightarrow$ architecture}

Le tableau ci-dessous synth\'{e}tise la correspondance entre chaque exigence
issue de la probl\'{e}matique et la couche architecturale qui y r\'{e}pond.

\begin{table}[h!]
\centering
\small
\begin{tabularx}{\textwidth}{l L{4cm} X}
\toprule
\textbf{R\'{e}f.} & \textbf{Exigence} & \textbf{R\'{e}ponse architecturale} \\
\midrule
E1 & D\'{e}tection temps-r\'{e}el,
     F1 $\geq$ 88\% &
     Couche D\'{e}tection ML : RF + XGBoost entra\^{i}n\'{e}s sur 11 features,
     10-fold CV, F1 = 99.92\%, FPR = 0.11\%. \\
E2 & Apprentissage comportemental &
     Couche Capture (agr\'{e}gation par fen\^{e}tre temporelle) +
     Couche D\'{e}tection (3 mod\`{e}les ML, pipeline de pr\'{e}traitement). \\
E3 & R\'{e}ponse autonome gradu\'{e}e &
     Couche R\'{e}ponse Active : cartographie de s\'{e}v\'{e}rit\'{e},
     4 modules de d\'{e}fense (Core Deception, MTD, Traffic Shaper, IP Manager). \\
E4 & Tableau de bord SOC temps-r\'{e}el &
     Couche Observabilit\'{e} : Flask + Socket.IO, 3 endpoints POST,
     dashboard D3.js, logs JSONL avec fallback local. \\
E5 & Explicabilit\'{e} des d\'{e}cisions &
     Couche D\'{e}tection : feature importance (Random Forest), scores de
     confiance affich\'{e}s dans le dashboard. \\
E6 & Isolation du laboratoire &
     Couche Capture : VirtualBox, r\'{e}seau 192.168.100.0/24 isol\'{e},
     mode offline pour d\'{e}mo sur Windows. \\
\bottomrule
\end{tabularx}
\caption{Tra\c{c}abilit\'{e} compl\`{e}te : chaque exigence issue de la
probl\'{e}matique trouve sa r\'{e}ponse dans une couche de l'architecture.}
\label{tab:tracabilite}
\end{table}

% ============================================================
\section{M\'{e}thodologie et cycle de vie}

L'architecture a \'{e}t\'{e} con\c{c}ue selon le cadre \textbf{CRISP-DM}
(\textit{Cross-Industry Standard Process for Data Mining}), qui structure le
projet en six phases :

\begin{enumerate}[nosep,leftmargin=1.5em]
  \item \textbf{Business Understanding} --- identification du besoin m\'{e}tier
        (r\'{e}duire la fatigue des analystes SOC), entretiens utilisateurs,
        d\'{e}finition des cibles quantitatives.
  \item \textbf{Data Understanding} --- exploration du dataset CICIDS2017,
        analyse de la distribution des classes, cartographie des features.
  \item \textbf{Data Preparation} --- pipeline de pr\'{e}traitement \`{a} 4
        \'{e}tapes avec contr\^{o}le anti-leakage.
  \item \textbf{Modeling} --- entra\^{i}nement et validation crois\'{e}e de
        RF, XGBoost et Isolation Forest.
  \item \textbf{Evaluation} --- m\'{e}triques sur l'ensemble de test,
        matrices de confusion, analyse du paradoxe IF.
  \item \textbf{Deployment} --- int\'{e}gration via le Bridge, dashboard SOC,
        d\'{e}fense active avec MTD et honeypot.
\end{enumerate}

Le d\'{e}tail complet est document\'{e} dans le rapport final AEGIS Entropy
(9 chapitres, Juin 2026).

\end{document}
